Small lesions in brain magnetic resonance imaging (MRI) are sparse,
spatially localized, and clinically important targets embedded within a
large volume of normal-appearing tissue. This results in an extreme
imbalance in voxel distribution, where lesion voxels constitute only a
small fraction of the image. Consequently, the learning process is
dominated by background and large structures, while signals from small
lesions are easily overlooked. Despite their size, small lesions play a
critical role in clinical assessment.

Deep learning methods, particularly convolutional neural networks (CNNs),
have substantially advanced medical image segmentation. Architectures
such as U-Net and its 3D extensions have become standard for volumetric
analysis~\cite{navab_u-net_2015}, while automated frameworks such as
nnU-Net have demonstrated strong performance across a wide range of
biomedical datasets through systematic pipeline
adaptation~\cite{isensee_nnu-net_2021}. More recently, transformer-based
and hybrid CNN--transformer models have further expanded the design space
for medical segmentation~\cite{hatamizadeh_unetr_2021,
hatamizadeh_swin_2022}. However, despite these advances, segmentation of
small pathological structures remains challenging even for
state-of-the-art models.

A primary difficulty arises from this imbalance. In brain MRI, lesion
voxels typically constitute only a very small fraction of the image
volume, causing standard training objectives to be dominated by
background regions. As a result, models tend to prioritize large
structures while under-representing small lesions during optimization.
To address this issue, several loss functions have been proposed. The
Tversky loss~\cite{salehi_tversky_2017} generalizes the Dice formulation
by introducing asymmetric weighting between false positives and false
negatives, allowing the objective to emphasize missed lesions and improve
recall. The Focal loss~\cite{lin_focal_2017} shifts training focus toward
hard examples by reducing the contribution of well-classified samples,
and its adaptation, the Focal Tversky
loss~\cite{abraham_novel_2019}, further enhances sensitivity and boundary
delineation in difficult regions. In addition, Yeung et
al.~\cite{yeung_calibrating_2023} propose DSC++, a modified Dice-based
objective that penalizes overconfident predictions, resulting in improved
calibration while preserving strong segmentation performance.

Although these methods improve voxel-level optimization under class
imbalance, they treat segmentation as a voxel-wise prediction problem and
do not explicitly model lesion structure. Small lesions typically appear
as discrete connected components scattered throughout the brain, and
their contribution to voxel-level objectives remains disproportionately
small. Consequently, models can achieve high overlap scores while still
missing clinically relevant small lesions, highlighting a fundamental
limitation of purely voxel-level supervision.

In the context of multiple sclerosis segmentation, several approaches
have explored architectural and ensemble strategies to improve
robustness. Wiltgen et al.~\cite{wiltgen_lst-ai_2024} proposed LST-AI,
an ensemble of 3D U-Net models trained on multimodal MRI with combined
losses, demonstrating improved generalization. Dereskewicz et
al.~\cite{dereskewicz_flames_2025} introduced FLAMeS, an nnU-Net-based
ensemble framework with strong performance on FLAIR MRI. Hashemi et
al.~\cite{hashemi_asymmetric_2019} proposed an asymmetric Tversky-based
loss within a 3D FC-DenseNet to improve the precision--recall trade-off
under severe imbalance. While effective, these approaches remain
fundamentally voxel-level and therefore do not explicitly address the
structural nature of small lesions.

More recent work has explored incorporating higher-level information into
the training objective. Region-based formulations modulate supervision
based on spatial context, such as adaptive region-level
losses~\cite{chen_adaptive_2023} and region-wise rectified
frameworks~\cite{valverde_region-wise_2023}. Zhang et
al.~\cite{zhang_all-net_2021} introduced lesion-aware modeling through
anatomical priors, while Javed et al.~\cite{aqib_javed_novel_2025}
proposed dual-coefficient regularization to better handle instance-level
imbalance. In parallel, multiple instance learning (MIL) provides a
set-level formulation where supervision is defined over groups of
instances rather than individual voxels. This has been explored in
medical imaging through object-aware MIL frameworks~\cite{liu_object_2026},
graph-based MIL approaches~\cite{huang_geometric_2026}, and weakly
supervised classification models~\cite{pal_deep_2021}. Although these
methods highlight the importance of object-level signals, they are not
designed to jointly optimize voxel-level segmentation and lesion-level
detection in a unified manner.

Overall, existing approaches either focus on voxel-level reweighting or
introduce partial higher-level supervision, but do not simultaneously
address structural imbalance and lesion-level detection. This suggests
that improving small lesion segmentation requires a unified objective
that integrates both component-level and instance-level information.

To address this gap, we propose a loss formulation specifically designed
for small lesion segmentation in brain MRI. First, we introduce a
Component-Adaptive Tversky (CAT) term that reweights the contribution of
individual lesion components, ensuring that small structures have a
stronger influence during training. Second, we incorporate a lesion-level
Multiple Instance Learning (MIL) term that encourages the model to detect
lesions at the component level. These two components are combined into a
unified CATMIL loss, which jointly optimizes voxel-level accuracy and
lesion-level detection.

We evaluate the proposed method on the MSLesSeg dataset for multiple
sclerosis lesion segmentation. Experimental results show that the
proposed loss improves the detection of small lesions while maintaining
competitive segmentation accuracy and stable error behavior compared to
standard objectives. These findings indicate that incorporating
structure-level and lesion-level supervision provides an effective and
practical strategy for improving small lesion segmentation.

The main contributions of this study are twofold. First, we propose a
loss formulation for small lesion segmentation that integrates a
component-adaptive term and a lesion-level supervision term within a
unified objective (CATMIL), enabling joint optimization of voxel-level
accuracy and lesion-level detection. Second, we provide a comprehensive
multi-level evaluation of segmentation behavior under severe class
imbalance, including voxel-level accuracy, lesion-level detection,
small-lesion recall, and detailed false positive and false negative
analysis. This evaluation reveals the strengths and trade-offs of
different objectives and demonstrates that the proposed formulation
improves small lesion detection while maintaining stable overall
performance.